% Antecedentes recuperados. Fragmento autónomo en sus definiciones;
% requiere amsmath, amssymb, amsthm, booktabs y longtable del documento principal.
\section{Realización finita de transporte con memoria de fase}
\label{iii:hist:antecedentes}

Los dos desarrollos siguientes reúnen mecanismos complementarios de la
construcción: una transferencia finita sobre posiciones APP, modos y fase,
y una colección de lectores cinemáticos y constitutivos sobre dos canales
angulares. Cada realización conserva su espacio de estados, sus operaciones
y sus parámetros. El operador finito tiene el dominio definido a
continuación; el TPK enriquecido del cuerpo principal conserva además sus
propias coordenadas y prolongaciones. Los lectores recíprocos se presentan
como realizaciones auxiliares con fórmulas y dominio declarados.

\subsection{Transferencia finita sobre posiciones, modos y fase}
\label{iii:hist:transferencia}

Sea $\operatorname{dr}_9(n)$ el representante en $\{1,\ldots,9\}$ de
$n\bmod9$, con $9$ como representante del residuo nulo. En las posiciones
$i,j\in\{0,\ldots,8\}$ se consideran las dos tablas
\[
 A_+(i,j)=\operatorname{dr}_9(i+j+2),\qquad
 A_\times(i,j)=\operatorname{dr}_9((i+1)(j+1)).
\]
Las entradas evaluadas aquí son enteros positivos. Por ello, la convención
histórica alternativa $\operatorname{dr}_9(0)=0$ no interviene en estas tablas.
El estado finito es
\[
 \Omega=(\mathbb Z/9\mathbb Z)^2\times\{+,\times\}
 \times\{N,E,S,O\}\times\mathbb Z/3\mathbb Z,
 \qquad |\Omega|=1944.
\]
Escribimos $x=(i,j,m,d,\tau)$. La transición conserva el orden siguiente:
se elige $d'$, se desplaza la posición, se lee APP con el modo anterior,
se reduce la lectura módulo $3$, se actualiza el modo y se avanza la fase.
Con
\[
 \delta_N=(-1,0),\quad\delta_E=(0,1),\quad
 \delta_S=(1,0),\quad\delta_O=(0,-1),
\]
las operaciones son
\begin{align*}
 (i',j')&=(i,j)+\delta_{d'}\pmod9,\\
 a&=A_m(i',j')\pmod3,\\
 m'&=\begin{cases}+,&a=1,\\\times,&a=2,\\m,&a=0,\end{cases}
 &\tau'&=\tau+1\pmod3.
\end{align*}
Así, $F_{d'}(x)=(i',j',m',d',\tau')$. El coste de giro y cambio de modo es
\[
 c_\lambda(x,d')=\mathbf1_{d'\ne d}+\lambda\mathbf1_{m'\ne m},
 \qquad \beta>0,\quad\lambda\geq0,
\]
y el operador actúa sobre la base de estados mediante
\begin{equation}
 T_{\beta,\lambda}e_x
 =\sum_{d'\in\{N,E,S,O\}}
 e^{-\beta c_\lambda(x,d')}e_{F_{d'}(x)}.
 \label{iii:hist:T}
\end{equation}
Los parámetros $\beta,\lambda$ son parámetros declarados de este modelo
finito. La tabla de referencia utiliza $(\beta,\lambda)=(1,1)$; el caso
$\lambda=0$ pertenece a las pruebas de sensibilidad.

\begin{proposition}[Equivarianza del antecedente finito]
\label{iii:hist:equivarianza}
La acción
\[
 \Phi_{u,v,w}(i,j,m,d,\tau)
 =(i+3u,j+3v,m,d,\tau+w),\qquad
 (u,v,w)\in(\mathbb Z/3\mathbb Z)^3,
\]
conmuta con $T_{\beta,\lambda}$.
\end{proposition}
\begin{proof}
Se cumple $\operatorname{dr}_9(n)\equiv n\pmod3$. Por tanto,
$A_+(i,j)\equiv i+j+2$ y
$A_\times(i,j)\equiv(i+1)(j+1)\pmod3$.
La traslación por múltiplos de tres preserva ambas lecturas reducidas.
También conmuta con el desplazamiento elegido y preserva el modo anterior,
el modo posterior y las dos direcciones. Conserva, en consecuencia, cada
indicador del coste. Finalmente, sumar $w$ a la fase conmuta con sumar una
unidad. Cada transición y su peso se transportan a la transición y al peso
correspondientes; la suma en \eqref{iii:hist:T} conserva esa igualdad.
\end{proof}

\begin{theorem}[Descomposición completa por caracteres]
\label{iii:hist:descomposicion}
Para $\omega=e^{2\pi i/3}$, el espacio $\mathbb C^\Omega$ se descompone
unitariamente en veintisiete subespacios de dimensión $72$. Sus bloques
satisfacen
\[
 T(k_a,k_b,k_t)=\omega^{k_t}T(k_a,k_b,0),
 \qquad k_a,k_b,k_t\in\{0,1,2\}.
\]
\end{theorem}
\begin{proof}
Escribimos $i=3a+r$, $j=3b+s$, con $a,b,r,s\in\{0,1,2\}$.
Cada fibra fina $u=(r,s,m,d)$ tiene $27$ estados $(a,b,\tau)$.
Definimos la isometría $U_k:\mathbb C^{72}\to\mathbb C^\Omega$ por
\[
 U_ke_u=\frac1{\sqrt{27}}
 \sum_{a,b,\tau=0}^2
 \omega^{-(k_aa+k_bb+k_t\tau)}e_{(3a+r,3b+s,m,d,\tau)}.
\]
La ortogonalidad de caracteres prueba $U_k^*U_\ell=\delta_{k\ell}I$ y
$\sum_kU_kU_k^*=I$. Si una transición cruza fronteras finas con incrementos
gruesos $\Delta a,\Delta b$, su coeficiente en $U_k^*TU_k$ es
\[
 e^{-\beta c_\lambda}\omega^{k_a\Delta a+k_b\Delta b+k_t}.
\]
Esto demuestra $TU_k=U_kT(k)$ y la descomposición del operador. El incremento
temporal es siempre uno, de modo que el último factor puede extraerse de
todo el bloque. Al elevar al cubo desaparece dicho factor temporal, aunque
los caracteres espaciales conservan su función.
\end{proof}

La descomposición anterior se apoya en la equivarianza y la ortogonalidad
de caracteres. La no negatividad del bloque trivial garantiza además un
autovalor real no negativo igual al radio espectral. Esta propiedad se
aplica al espacio completo, conservando también sus estados transitorios.

\subsection{Descomposición espectral y sensibilidad paramétrica}
\label{iii:hist:espectros}

Sean $\rho_1\geq\rho_2$ los dos módulos mayores, contados con
multiplicidad, y $g=\rho_1-\rho_2$. La siguiente tabla contiene los
nueve bloques espaciales; el teorema anterior determina los veintisiete
bloques temporales completos: para cada fila y cada $t=0,1,2$, el autovalor
dominante es $\omega^t\mu_1(k_a,k_b,0)$, mientras $\rho_1,\rho_2,g$
permanecen iguales. Esta forma conserva explícitamente las tres fases.
En esta primera tabla se fija $\beta=\lambda=1$.

\begin{center}
\small
\begin{tabular}{ccrrrr}
\toprule
$k_a$&$k_b$&$\mu_1(k_a,k_b,0)$&$\rho_1$&$\rho_2$&$g$\\
\midrule
0&0&1.730888896&1.730888896&1.083076285&0.647812612\\
0&1&1.459962951&1.459962951&1.116843785&0.343119167\\
0&2&1.459962951&1.459962951&1.116843785&0.343119167\\
1&0&1.459962951&1.459962951&1.116843785&0.343119167\\
1&1&-1.508515719&1.508515719&1.001353790&0.507161929\\
1&2&-1.514010372&1.514010372&0.971977167&0.542033205\\
2&0&1.459962951&1.459962951&1.116843785&0.343119167\\
2&1&-1.514010372&1.514010372&0.971977167&0.542033205\\
2&2&-1.508515719&1.508515719&1.001353790&0.507161929\\
\bottomrule
\end{tabular}
\end{center}

Los signos de los autovalores dominantes y sus fases forman parte de la
salida: los cuatro sectores con $k_a,k_b\ne0$ tienen un líder real negativo
para $k_t=0$. Su módulo conserva la tasa espectral y su signo conserva una
información adicional que interviene al componer las iteraciones.

Los cinco mapas completos de sensibilidad, con filas $k_a$ y columnas
$k_b$, son los siguientes; todos corresponden a $k_t=0$:
\begin{align*}
\rho(1,1)&=\begin{pmatrix}
1.730888896&1.459962951&1.459962951\\
1.459962951&1.508515719&1.514010372\\
1.459962951&1.514010372&1.508515719
\end{pmatrix},\\
\rho(1,0)&=\begin{pmatrix}
2.103638324&1.685504327&1.685504327\\
1.685504327&1.853245966&1.853245966\\
1.685504327&1.853245966&1.853245966
\end{pmatrix},\\
\rho(1,2)&=\begin{pmatrix}
1.656915306&1.433869155&1.433869155\\
1.433869155&1.423885161&1.432116157\\
1.433869155&1.432116157&1.423885161
\end{pmatrix},\\
\rho(0.5,1)&=\begin{pmatrix}
2.480049445&2.124801655&2.124801655\\
2.124801655&2.281861809&2.286772914\\
2.124801655&2.286772914&2.281861809
\end{pmatrix},\\
\rho(2,1)&=\begin{pmatrix}
1.226049377&1.131390710&1.131390710\\
1.131390710&1.028041508&1.028201900\\
1.131390710&1.028201900&1.028041508
\end{pmatrix}.
\end{align*}
Aquí $\rho(\beta,\lambda)$ denota el mapa de radios, no el operador de
transferencia. El intercambio $i\leftrightarrow j$ junto con el intercambio
correspondiente de direcciones preserva APP, las transiciones y el coste;
por ello estos mapas son simétricos. La variación de los parámetros modifica los radios manteniendo
la descomposición por caracteres, que depende de la simetría de las
operaciones. Las tablas muestran aproximaciones numéricas de espectros
finitos. Las identidades de transición, equivarianza y factorización se
sostienen por las pruebas anteriores; sus consecuencias numéricas se
conservan junto con los datos completos de las veintisiete componentes.

\subsection{Lectores recíprocos y pruebas de sensibilidad}
\label{iii:hist:n6d}

Para reproducir la realización angular antecedente se conservan las medidas
numéricas decimales en grados
\[
 a_{\mathrm h}=7.297352569283802,\qquad
 c_{\mathrm h}=2.123738338968646.
\]
Los ángulos y la coordenada adimensional correspondiente se distinguen mediante
\[
 A_{\mathrm h}=a_{\mathrm h}{}^\circ,\qquad
 C_{\mathrm h}=c_{\mathrm h}{}^\circ,\qquad
 \alpha_{\mathrm h}=a_{\mathrm h}/1000.
\]
Estos datos permiten cotejar el antecedente; las constantes del desarrollo
vigente conservan su generación anterior APP--TRIT--TPK y su representación
posterior. La reproducción histórica no reemplaza esa generación por los
decimales escritos aquí. La conversión angular y los canales utilizados son
\[
 x_{\mathrm h}=\frac{\pi a_{\mathrm h}}{180},\qquad
 y_{\mathrm h}=\frac{\pi c_{\mathrm h}}{180},\qquad
 q_\pm=e^{-(x_{\mathrm h}\pm y_{\mathrm h})},
 \qquad0<q_+<q_-<1.
\]
El primer lector expresa
\[
 \epsilon_{90}=q_-^{90}-q_+^{90},\qquad
 c_\pm=\frac{c_{\mathrm{ref}}}{1\pm\epsilon_{90}}.
\]
Aquí $c_{\mathrm{ref}}$ sólo normaliza la prueba cinemática. Para reproducir
sus unidades se utilizó el mismo valor de referencia histórico
$299792458\,\mathrm{m\,s^{-1}}$.

\begin{proposition}[Media armónica y producto de dos hojas]
\label{iii:hist:lectores}
Si $|\epsilon_{90}|<1$, la media armónica de $c_+$ y $c_-$ es
$c_{\mathrm{ref}}$. En el lector constitutivo histórico
\[
 \mu_{\mathrm h}^{\pm}=e^{\pm2\alpha_{\mathrm h}},\qquad
 \varepsilon_{\mathrm h}^{\pm}=e^{\mp2\alpha_{\mathrm h}},
\]
se tiene $\mu_{\mathrm h}^{\pm}\varepsilon_{\mathrm h}^{\pm}=1$ y
$Z_+/Z_-=e^{4\alpha_{\mathrm h}}$.
\end{proposition}
\begin{proof}
La suma $c_+^{-1}+c_-^{-1}=2/c_{\mathrm{ref}}$ elimina el término orientado.
En cada hoja, la suma de los dos exponentes constitutivos es cero.
Por otra parte, $Z_\pm=\sqrt{\mu_{\mathrm h}^{\pm}/
\varepsilon_{\mathrm h}^{\pm}}=e^{\pm2\alpha_{\mathrm h}}$.
El cociente de impedancias conserva, por tanto, la orientación que elimina
el producto constitutivo.
\end{proof}

El segundo lector histórico procede de series geométricas convergentes:
\begin{align*}
 S_a(q)&=\frac{q^a}{1-q^{3a}}
       =\sum_{j\geq0}q^{a(1+3j)},\\
 \delta_{a,b}(q)&=S_a(q)-S_b(q),\\
 R_{a,b}&=\frac{\delta_{a,b}(q_-)-\delta_{a,b}(q_+)}
 {q_-^a-q_+^a}.
\end{align*}
La convergencia vale para $|q|<1$; el denominador del último cociente es
positivo para los canales y los enteros positivos considerados. Los lectores
auxiliares históricos son
\[
 \zeta_{12}=1-R_{90,120}^2,\qquad
 \Delta_{4,\mathrm h}=\frac{x_{\mathrm h}-y_{\mathrm h}}{810}
 \left(1-\frac{7\pi}{12\cdot729}\right),\qquad
 \xi_{\mathrm h}=1-\frac{27}{160}\Delta_{4,\mathrm h}.
\]
El subíndice $\mathrm h$ mantiene explícita esta definición histórica de
torsión, evitando identificarla por el símbolo con otro lector vigente.

La evaluación a partir de las cadenas decimales declaradas da
\begin{align*}
 q_+&=0.848377942576404356\ldots,&
 q_-&=0.913660151150557274\ldots,\\
 \epsilon_{90}&=0.000295169439289827959\ldots,&
 R_{90,120}&=0.933314535237713061\ldots,\\
 \zeta_{12}&=0.128923978314011665\ldots,&
 \xi_{\mathrm h}&=0.999981235497802379\ldots.
\end{align*}
La tabla completa de perturbaciones discretas es
\begin{center}
\begin{tabular}{rrr}
\toprule
$a$&$b$&$R_{a,b}$\\\midrule
90&120&0.933314535238\\
90&119&0.927013603649\\
90&121&0.939071551912\\
89&120&0.939066204789\\
91&120&0.927019444194\\
30&40&0.569215041549\\
60&80&0.834174862952\\
\bottomrule
\end{tabular}
\end{center}
La sensibilidad angular distingue las dos coordenadas del par:
\begin{center}
\begin{tabular}{lrrr}
\toprule
Variable&Variación relativa&$R_{90,120}$&Variación en ppm\\\midrule
$A_{\mathrm h}$&$-0.001$&0.933059250359&$-273.525022305$\\
$A_{\mathrm h}$&$+0.001$&0.933568846564&$+272.481908792$\\
$C_{\mathrm h}$&$-0.001$&0.933388163678&$+78.889202964$\\
$C_{\mathrm h}$&$+0.001$&0.933240822366&$-78.979667757$\\
\bottomrule
\end{tabular}
\end{center}

Las evaluaciones se realizaron con $80$ y $120$ cifras de precisión a partir
de las cadenas decimales declaradas. La precisión de evaluación conserva
la especificación finita de aquellos datos angulares.

Finalmente, estos lectores mantienen su lugar en la genealogía documental.
El operador constitutivo vigente
$V=(I-T^{90})(I-T^{120})^{-1}$ evalúa
$(1-q_\pm^{90})/(1-q_\pm^{120})$ en sus dos hojas. El funcional orientado
vigente utiliza $12S_{90}-S_{120}$. La realización antecedente conserva el lector
$S_{90}-S_{120}$ y su normalización por $q_-^{90}-q_+^{90}$.
La relación entre ellos se expresa manteniendo cada operación y sus pesos,
en lugar de sustituirlos por un único número o un mismo símbolo.
